\chapter{סקירת ספרות ומצב קיים}\label{ch:lit}

הסקירה מסכמת ארבעה גופי ידע: הטכנולוגיה לזיהוי תמונות סינתטיות, תשתית המקור (\en{Provenance}), הכלכלה של תוכן מזויף, וכללי האתיקה העיתונאית. בסופה מוגדר הפער שאף פתרון קיים אינו סוגר.

\section{גישות טכנולוגיות לאימות תמונה}
בספרות ובתעשייה מקובלות שלוש משפחות של פתרונות, ולכל אחת חולשה מובנית:

\subsection{זיהוי פסיבי (גלאים סטטיסטיים)}
מסווגים, לרוב רשתות עצביות, שמחפשים ״טביעות״ שמחולל משאיר בתמונה: דפוסים בתחום התדר, אי עקביות בתאורה ובצללים, או מבנה רעש שאינו מתאים לחיישן מצלמה. היתרון הוא שאין צורך בשיתוף פעולה של יוצר התמונה. החיסרון חמור: גלאי שאומן על מחוללים מסוימים מאבד דיוק מול מחולל חדש, ועיבוד פשוט כמו דחיסה, הקטנה או צילום מסך מוחק חלק מהטביעות. מחקר משנת \num{2025} בדק גלאים מתקדמים על מאגר החלפות פנים ״מהשטח״ ומצא שביצועיהם מתקרבים לניחוש אקראי. עוד נמצא שעיבוד של הגדלת רזולוציה פוגע בהם משמעותית~\cite{ren2025}.

\textbf{השלכה לפרויקט:} אסור לבנות מוצר שמבטיח ״זיהוי בינה מלאכותית״ על סמך גלאי יחיד. גלאי הוא אות אחד מתוך כמה, וחייב להיות מכויל ולהציג אי ודאות.

\subsection{סימני מים מוטמעים}
המחולל עצמו מטמיע בתמונה סימן בלתי נראה שמזהה אותה כסינתטית. חוק הבינה המלאכותית האירופי מחייב ספקי מחוללים לסמן פלטים ״בפורמט קריא למכונה״ שנשאר ניתן לזיהוי, ומחוללים שכבר היו בשוק לפני \num{2} באוגוסט \num{2026} קיבלו ארכה עד \num{2} בדצמבר \num{2026}~\cite{euaiact2026}. החולשה: הסימון תלוי ברצון הספק, ומחוללים בקוד פתוח שרצים על מחשב פרטי אינם מחויבים בפועל.

\subsection{תעודות מקור (\en{C2PA})}
גישה הפוכה: במקום לחפש סימן של זיוף, חותמים קריפטוגרפית על האמת. המצלמה או התוכנה מוסיפה לקובץ ״תעודת מקור״ שמתעדת מי יצר אותו, מתי ובאילו כלים נערך~\cite{c2paspec}. ב-\num{2026} קיימת חתימה במצלמות של \en{Leica}, \en{Sony}, \en{Canon}, \en{Fujifilm} ו-\en{Panasonic}, ו-\en{Google Pixel 10} חותם כברירת מחדל על כל תמונה~\cite{c2pa2026}. \en{Canon} השיקה במאי \num{2026} מערכת ייעודית למערכות חדשות שנבדקה עם \en{Reuters}~\cite{canon2026}.

שלוש מגבלות מתועדות: (\num{1}) רשתות חברתיות מסירות לעתים קרובות את המטא דאטה בהעלאה ובדחיסה; (\num{2}) פחות מ-\num{1\%} מתמונות החדשות בעולם נושאות היום תעודה; (\num{3}) התעודה מתעדת רק את מה שהוצהר בעת החתימה, ואינה מזהה תוכן סינתטי שלא נחתם~\cite{c2pa2026}. נוסף על כך, החתימה במצלמות \en{Nikon} הושעתה והתעודות שלה בוטלו~\cite{c2pa2026}, כלומר גם שרשרת האמון עצמה פגיעה (טבלה~\ref{tab:approaches}).

\begin{table}[!htb]\centering
\caption{לכל גישה טכנולוגית חולשה מובנית, ולכן אין להסתמך על אחת בלבד}\label{tab:approaches}
\begin{tabularx}{\linewidth}{@{}l Y Y Y@{}}
\toprule
\textbf{קריטריון} & \textbf{זיהוי פסיבי} & \textbf{סימן מים} & \textbf{תעודת מקור \en{C2PA}} \\
\midrule
מה מוכח & הסתברות לזיוף & שהתמונה סינתטית & שהתמונה אמיתית ומה נערך בה \\
תלות בשיתוף פעולה & אין & ספק המחולל & המצלמה והתוכנה \\
עמידות לדחיסה ולצילום מסך & נמוכה עד בינונית & בינונית & נמוכה (המטא דאטה נמחק) \\
כיסוי תמונות קיימות & כמעט מלא & רק ממחוללים מסמנים & פחות מ-\num{1\%} מתמונות החדשות~\cite{c2pa2026} \\
סיכון עיקרי & התרעות שווא ופספוסים & מחוללים לא מסמנים & היעדר תעודה אינו הוכחה לזיוף \\
\bottomrule
\end{tabularx}
\end{table}

\textbf{מסקנה:} שלוש הגישות משלימות זו את זו. תעודת מקור נותנת ודאות כשהיא קיימת. סימן מים נותן ודאות כשהוא נמצא. הגלאי מכסה את כל השאר, אבל בהסתברות בלבד. ארכיטקטורה נכונה מריצה את שלושתן ומשקללת אותן, ולכן נבחרה בפרק~\ref{ch:alt}.

\section{הכלכלה של תמונות מזויפות}
\en{DiResta} ו-\en{Goldstein} בדקו \num{125} עמודי פייסבוק שכל אחד מהם פרסם לפחות \num{50} תמונות בינה מלאכותית. לעמודים היו בממוצע \num{146,681} עוקבים, והתמונות צברו יחד מאות מיליוני חשיפות. פוסט אחד הגיע ל-\num{40} מיליון צפיות ו-\num{1.9} מיליון תגובות~\cite{diresta2024}. החוקרים זיהו שני מודלים: ״ספאמרים״ שמפנים קוראים לאתרי תוכן דל כדי להרוויח מפרסום, ו״רמאים״ שמוכרים מוצרים שאינם קיימים. רוב התמונות לא תויגו כסינתטיות.

ברמת האתרים, \en{NewsGuard} זיהתה נכון למרץ \num{2026} יותר מ-\num{3,000} ״חוות תוכן״ אוטומטיות, ומאות נוספות בכל חודש~\cite{newsguard2026}. אתרים כאלה ממומנים בעיקר בפרסום פרוגרמטי, כלומר במודעות של מותגים שלרוב אינם יודעים איפה המודעה שלהם מופיעה. כאן נמצא מנוף מרכזי: אם המפרסמים יפסיקו לממן את הטראפיק המזויף, התמריץ נחלש. זה הבסיס למוצר ״ציון שלמות״ לרשתות פרסום (פרק~\ref{ch:business}).

\section{אמון הציבור}
סקר של מכון רויטרס ואוניברסיטת אוקספורד משנת \num{2025} בשש מדינות מצא ש-\num{12\%} בלבד מרגישים בנוח עם חדשות שהופקו כולן בבינה מלאכותית, לעומת \num{62\%} שמעדיפים תוכן שהופק בידי אדם. הנוחות עולה ל-\num{43\%} כשאדם מוביל את התהליך ובינה מלאכותית רק מסייע. הציבור מקבל בקלות שימושים ״מאחורי הקלעים״, כמו תרגום, יותר משימושים גלויים כמו תמונות שנוצרו בבינה מלאכותית~\cite{reuters2025}. מכאן עיקרון תכנון: העורך האנושי נשאר מקבל ההחלטה, והמערכת מסייעת לו ומתעדת.

\section{כללי אתיקה קיימים}

ראו טבלה~\ref{tab:ethics}.

\begin{table}[!htb]\centering
\caption{מה קובעים כללי האתיקה הקיימים לגבי תמונות ובינה מלאכותית}\label{tab:ethics}
\begin{tabularx}{\linewidth}{@{}l Y Y@{}}
\toprule
\textbf{מקור} & \textbf{מה נקבע} & \textbf{מה חסר} \\
\midrule
קוד האתיקה של \en{SPJ} (\num{2014})~\cite{spj2014} & ״לעולם אל תעוות במכוון עובדות או הקשר, כולל מידע חזותי״; ״סמן בבירור איורים ושחזורים״ & לא נכתב לעידן הבינה המלאכותית; אין נוהל לאימות \\
\en{AP} (\num{2023})~\cite{ap2023} & אסור לשנות צילום, וידאו או אודיו בכלים גנרטיביים; איור בינה מלאכותית חייב תיוג מתאים & מדיניות של ארגון אחד; אין כלי אכיפה \\
אמנת פריז (\num{2023})~\cite{paris2023} & \num{10} עקרונות, ובהם: האדם מקבל ההחלטה, גילוי שימוש משמעותי בבינה מלאכותית, עקיבות מקור, הבחנה ברורה בין אותנטי לסינתטי & עקרונות ולא נהלים; אין הסמכה \\
חוק הבינה המלאכותית האירופי, סעיף \num{50}~\cite{euaiact2026,euaiact} & חובת גילוי ל״דיפ פייק״; לפי טיוטת ההנחיות של הנציבות: מבחן של ארבעה רכיבים, בלי צורך בכוונה להטעות & לא מגדיר איך מאמתים בפועל \\
ישראל~\cite{mako2024} & אין חקיקה ייעודית לבינה מלאכותית; לפי עמדת הרשות להגנת הפרטיות, ״דיפ פייק״ שעלול להיתפס כאותנטי כפוף לדיני הפרטיות & אין חובת תיוג ואין תקן מקצועי ייעודי \\
\bottomrule
\end{tabularx}
\end{table}

הדפוס חוזר: יש הסכמה רחבה על העקרונות (אמת, שקיפות, אחריות אנושית), ואין כמעט תרגום שלהם לנוהל עבודה שאפשר לבצע ולבקר. התקן \en{TRUTH-IMG} נועד לסגור את הפער הזה.

\section{ניתוח אתי: למה ״כולם עושים את זה״ אינו נימוק}
פרסום תמונה סינתטית כדי לייצר טראפיק נבחן כאן לפי שלוש מסגרות אתיות קלאסיות. כל אחת מגיעה לאותה מסקנה מכיוון אחר, וזה מה שמצדיק כלל מחייב ולא רק המלצה (טבלה~\ref{tab:frameworks}).

\begin{table}[!htb]\centering
\caption{שלוש מסגרות אתיות ומה הן קובעות על תמונה סינתטית שמוצגת כתיעוד}\label{tab:frameworks}
\small
\begin{tabularx}{\linewidth}{@{}l Y Y@{}}
\toprule
\textbf{מסגרת} & \textbf{השאלה שהיא שואלת} & \textbf{המסקנה לגבי ״תמונה למען טראפיק״} \\
\midrule
דאונטולוגית (חובה) & האם המעשה מכבד את הקורא כאדם בעל זכות לאמת? & לא. הקורא משמש אמצעי להכנסה. ההטעיה פסולה גם אם לא נגרם נזק נראה לעין. \\
תוצאתית (תועלת) & האם סך התוצאות חיובי? & לא. הרווח בטראפיק קצר וממוקד במפרסם, הנזק לאמון מצטבר ומפוזר על כל הציבור, וגם ״דיבידנד השקרן״ מתחזק. \\
מידות טובות (אופי) & מה היה עושה עיתונאי ישר ואמיץ? & היה מעכב את הפרסום, מודה בספק ומעדיף להיות צודק ולא ראשון. \\
\bottomrule
\end{tabularx}
\end{table}

שלושה טיעונים נגדיים נפוצים, והמענה להם:
\begin{enumerate}
\item \textbf{״זו רק תמונה להמחשה.״} אם כך, יש לתייג אותה כאיור (קטגוריה ד׳). הבעיה אינה בשימוש בתמונה סינתטית, אלא בהצגתה כתיעוד. התיוג מבטל את ההטעיה ומשאיר את ההמחשה.
\item \textbf{״המתחרים כבר פרסמו.״} זה טיעון של דילמת אסיר: לכל מערכת משתלם לפרסם, ולכולן יחד משתלם שאף אחת לא תפרסם. דילמה כזו נפתרת בכלל משותף שכולם מחויבים לו ובמנגנון שמתגמל את מי שמקיים אותו. זה בדיוק התפקיד של התקן ושל ״ציון השלמות״ למפרסמים.
\item \textbf{״הקוראים יודעים שזה בינה מלאכותית.״} הנתונים מראים שלא: רוב התמונות הוויראליות שנבדקו לא תויגו כלל~\cite{diresta2024}, ותמונת הוריקן ״הילדה עם הגור״ הופצה ושימשה למסר פוליטי בלי שום סימון שהיא סינתטית~\cite{verify2024}.
\end{enumerate}

\section{פתרונות קיימים בשוק}
נבדקו המתחרים הישירים והעקיפים כפי שהם מציגים את עצמם בשוק בנובמבר \num{2026}. הממצא החשוב ביותר: יש כבר שירות אחד שבנוי בדיוק למערכות חדשות, ולכן היתרון של \en{VeriLens} לא יכול להיות ״זיהוי טוב יותר״ (טבלה~\ref{tab:rivals}).

\begin{table}[!htb]\centering
\caption{המתחרים: במה הם מתמקדים ומה ידוע על מחיר וגיוסים}\label{tab:rivals}
\small
\begin{tabularx}{\linewidth}{@{}l Y Y l@{}}
\toprule
\colhead{מתחרה} & \colhead{מיקוד} & \colhead{מחיר או גיוס} & \colhead{מקור} \\
\midrule
\en{ImageWhisperer} & אימות תמונות למערכות חדשות: \num{18} מודלי זיהוי במקביל, חיפוש הפוך, אימות מיקום, מאגר זיופים מוכרים; לקוחות: \en{Axel Springer}, \en{Mediahuis NRC} & \num{500} עד \num{900} יורו בחודש & \cite{imagewhisperer} \\
\en{Reality Defender} & ארגונים, בנקים ומוקדי שירות; תמונה, וידאו, אודיו וטקסט & גיוס של כ-\num{33} מיליון דולר & \cite{techinsider2026} \\
\en{Hive} & סינון תוכן בפלטפורמות גדולות; הזיהוי הוא חלק ממוצר רחב & גיוס של כ-\num{120} מיליון דולר & \cite{techinsider2026} \\
\en{Sensity} & חקירות הונאה ואכיפת חוק, עם תיעוד ראייתי & גיוס של כ-\num{3.2} מיליון דולר & \cite{techinsider2026} \\
\en{GetReal Security} & הגנה על ארגונים מפני מדיה מזויפת & גיוס של \num{17.5} מיליון דולר & \cite{getreal2025} \\
\en{TrueMedia} & כלי חינמי לעיתונאים; נסגר בינואר \num{2025} והועבר לאוניברסיטת ג׳ורג׳טאון & חינם & \cite{truemedia2025} \\
\bottomrule
\end{tabularx}
\end{table}

לפי סקירה שמשווה את \en{Reality Defender}, \en{Hive} ו-\en{Sensity}, לאף אחת מהן אין מקרי שימוש מפורסמים במערכות חדשות~\cite{techinsider2026}. בישראל, \en{Clarity}, שעסקה בזיהוי זיופים, נרכשה באוגוסט \num{2026} ועברה לאימות זהות של עובדים~\cite{clarity2026}. בדיקה של \en{Columbia Journalism Review} מזכירה שגם הכלים הטובים הם ״שכבת אישור אחת״ בתהליך אימות רחב יותר, ולא תחליף לו~\cite{cjr2025}.

\begin{table}[!htb]\centering
\caption{מה מכסה כל פתרון: הזיהוי כבר קיים, התיעוד והתקן האתי חסרים}\label{tab:compet}
\small
\begin{tabularx}{\linewidth}{@{}Y c c c c c c@{}}
\toprule
\colhead{פתרון} & \colhead{\en{C2PA}} & \colhead{גלאים} & \colhead{הקשר} & \colhead{שילוב ב-\en{CMS}} & \colhead{יומן החלטות} & \colhead{תקן אתי} \\
\midrule
\en{ImageWhisperer} & לא צוין & \checkmark & \checkmark & \checkmark & חלקי & לא \\
ממשקי זיהוי כלליים (\en{Reality Defender}, \en{Hive}, \en{Sensity}) & לא & \checkmark & לא & חלקי & חלקי & לא \\
תוספי אימות לעיתונאים (\en{InVID-WeVerify}) & חלקי & חלקי & \checkmark & לא & לא & לא \\
בודקי תעודות מקור (\en{Content Credentials Verify}) & \checkmark & לא & לא & לא & לא & לא \\
מדיניות פנימית של מערכת חדשות (\en{AP}) & לא & לא & לא & לא & לא & חלקי \\
\textbf{\en{VeriLens} (מוצע)} & \checkmark & \checkmark & \checkmark & \checkmark & \checkmark & \checkmark \\
\bottomrule
\end{tabularx}
\end{table}

\section{הפער}
\begin{keybox}[הפער]
זיהוי תמונות למערכות חדשות כבר קיים בשוק, ולקוחות גדולים כבר משלמים עליו. מה שחסר הוא השכבה שמעל הזיהוי: אימות תעודות מקור כחלק מאותו תהליך, יומן החלטות אנושיות שמגן על מערכת החדשות, תיוג גלוי לקורא, ותקן אתי עם הסמכה. זו התרומה החדשה של \en{VeriLens}, ובנוסף שוק העברית והערבית, שבו אין היום ספק מקומי.
\end{keybox}
